% intro.tex -- funnel, pipeline, objective equation, open question
% (skills/paper-playbook/sections/03..06).
\section{Introduction}\label{sec:intro}
<<Field sentence with founding citations.>> <<Schemes of choice and why.>>
<<Why the expensive component exists.>> <<Lineage sentence citing the line at once.>>
<<Bottleneck sentence, ideally with a citable measurement.>>
Concretely, <<procedure>> factors into three stages:
\begin{itemize}
\item \textbf{Stage 1: <<name>>.} <<map>>, cost $\bigO(<<\cdot>>)$.
\item \textbf{Stage 2: <<name>>.} <<map>>, cost $\bigO(<<\cdot>>)$.
\item \textbf{Stage 3: <<name>>.} <<map>>, cost $\bigO(<<\cdot>>)$.
\end{itemize}
In this work we concentrate on Stage <<j>>. Every previous construction can be cast as
\begin{equation}\label{eq:intro:objective}
  \min_{P}\ \cost(P) \quad \text{s.t.}\quad <<functional specification>>.
\end{equation}
Whether <<further structure>> can lower~\eqref{eq:intro:objective} further has remained open.

\subsection{Related Work}\label{sec:related}
We overview prior optimisations for <<task>>, focusing on <<bottleneck>>. Accelerating
\eqref{eq:intro:objective} follows two paths: <<axis 1>> and <<axis 2>>.
\Cref{tab:related} summarises the asymptotic comparison.
% related_table.tex -- comparison table; "Ours" row must equal the abstract's formula.
\begin{table}[t]
\centering\scriptsize
\caption{Asymptotic comparison of <<methods>>. <<Define every symbol.>>
$^\dagger$<<scope restriction>>. $^\ddagger$<<scope restriction>>.}
\label{tab:related}
\renewcommand{\arraystretch}{0.92}\setlength{\tabcolsep}{3pt}
\begin{tabular}{llccc}
\toprule
Technique & Method & <<Cost 1>> & <<Cost 2>> & <<Coordinate>> \\
\midrule
\multirow{2}{*}{<<Axis 1>>} & <<[A]>> & $\bigO(\cdot)$ & $\bigO(\cdot)$ & <<(1,$\cdot$)>> \\
                            & <<[B]>>$^\dagger$ & $\bigO(\cdot)$ & $\bigO(\cdot)$ & <<($\cdot$,1)>> \\
\midrule
<<Axis 2>> & <<[C]>>$^\ddagger$ & $\bigO(\cdot)$ & $\bigO(\cdot)$ & <<($\cdot$,$\cdot$)>> \\
\midrule
\textbf{Ours} & & $\bigO(\cdot)$ & $\bigO(\cdot)$ & <<interior>> \\
\bottomrule
\end{tabular}
\end{table}


\noindent\textbf{<<Axis 1>>.} <<One sentence per work: authors, what, cost. Shared limitation.>>

\noindent\textbf{<<Recent lines>>.} Recently, two independent lines of work exploit <<structure>>:
\begin{itemize}
\item \textbf{<<Line A>>} restricts \emph{where} <<...>>; however, <<factual open point>>.
\item \textbf{<<Line B>>} restricts \emph{how} <<...>>, assuming <<...>>.
\end{itemize}

\noindent\textbf{Insight from <<structure>>.} <<Why the lines are compatible; where each
prior work sits in our framework; back to \Cref{tab:related}.>>

\noindent\textbf{Our work.} <<Framework sentence; three technical points; one results sentence.>>

\subsection{Our Contributions}\label{sec:contrib}
Inspired by <<observation>>, we determine <<characterisation>>. Building on it, we propose
<<Name>>, in which every prior <<object>> is a special case. Our main contributions are:
\begin{itemize}
\item \textbf{From <<structure>> to <<construction>>.} <<... coverage with one number (\Cref{tab:main}).>>
\item \textbf{A unified <<framework>>.} <<... placement of prior work ... selection rule.>>
\item \textbf{A <<algorithm>> with <<property>>.} <<... cost formula ...>>
\item \textbf{<<Metric>> improvements.} <<two numbers, verified in \Cref{tab:main}.>>
\end{itemize}

\subsection{Technique Overview}\label{sec:overview}
We begin by recalling the routes of \Cref{sec:related} in the notation of \Cref{sec:prelim}.

\noindent\textbf{<<Route 1>> (per paper).} <<formula, cost; what it leaves open.>>

\noindent\textbf{<<Route 2>> (per paper).} <<formula, cost; what it leaves open.>>

Taken together, <<...>>. Whether they compose therefore reduces to two questions:
(i)~<<...>>; (ii)~<<...>>. Insight~1 resolves (i) and Insight~2 resolves (ii),
yielding the map in \Cref{fig:overview}.

\begin{figure}[t]
\centering
\fbox{\parbox{0.8\linewidth}{\centering <<figure: prior works as corners, ours as the new region>>}}
\caption{<<One sentence: prior works are ...; our work occupies ...>>}
\label{fig:overview}
\end{figure}

\noindent\textbf{Insight 1: <<A complete declarative sentence.>>} <<...>>

\noindent\textbf{Insight 2: <<A complete declarative sentence.>>} <<...>>

\noindent\textbf{Putting it together.} <<Combined identity, literally as in \Cref{sec:construction}.>>
